\section{Danh sách User Story và tiêu chí chấp nhận}

Mỗi User Story dưới đây tương ứng với một task trong Product Backlog của Assignment 1. Phát biểu nêu đối tượng, nhu cầu và lợi ích; các tiêu chí chấp nhận (AC) mô tả kết quả có thể kiểm tra khi nghiệm thu. Những chức năng mở rộng của BrewLite được nêu tại story liên quan để phản ánh luồng hiện tại.

\subsection{US-01 -- Khởi tạo nền tảng (Task 1)}
\textbf{User Story}\par
Là một thành viên nhóm phát triển, tôi muốn có cấu trúc dự án và hướng dẫn chạy thống nhất, để cả nhóm có thể phát triển, tích hợp và kiểm thử trên máy cá nhân.

\textbf{Tiêu chí chấp nhận (Acceptance Criteria):}
\begin{enumerate}
  \item Repository có hai ứng dụng \texttt{frontend} (Next.js) và \texttt{backend} (NestJS), được quản lý bằng Git.
  \item Có \texttt{README.md} mô tả cấu trúc, các biến môi trường mẫu và cách khởi chạy các dịch vụ.
  \item Sau khi làm theo hướng dẫn, frontend hiển thị trang web và backend trả lời một API kiểm tra được.
\end{enumerate}

\subsection{US-02 -- API danh sách sản phẩm (Task 2)}
\textbf{User Story}\par
Là một khách hàng, tôi muốn nhận danh sách đồ uống đang bán, để biết những món có thể lựa chọn.

\textbf{Tiêu chí chấp nhận (Acceptance Criteria):}
\begin{enumerate}
  \item \texttt{GET /api/products} trả về JSON từ cơ sở dữ liệu, gồm mã, tên, ảnh và giá theo size của từng sản phẩm.
  \item Sản phẩm ngừng kinh doanh không xuất hiện trong menu dành cho khách.
  \item Khi không có sản phẩm phù hợp, API trả danh sách rỗng hợp lệ, không gây lỗi máy chủ.
\end{enumerate}

\subsection{US-03 -- Xem menu (Task 3)}
\textbf{User Story}\par
Là một khách hàng, tôi muốn xem menu đồ uống trên web, để tìm và chọn món phù hợp trước khi đặt.

\textbf{Tiêu chí chấp nhận (Acceptance Criteria):}
\begin{enumerate}
  \item Trang Menu gọi API sản phẩm và hiển thị các món dưới dạng lưới, có tên, ảnh và giá để nhận biết.
  \item Có trạng thái đang tải, danh sách rỗng và thông báo khi API không phản hồi.
  \item Chọn một món mở được trang chi tiết của đúng sản phẩm đó.
\end{enumerate}

\subsection{US-04 -- Tùy chỉnh đồ uống (Task 4)}
\textbf{User Story}\par
Là một khách hàng, tôi muốn chọn size và topping hợp lệ cho đồ uống, để món đặt phù hợp khẩu vị và tôi biết giá trước khi thêm vào giỏ.

\textbf{Tiêu chí chấp nhận (Acceptance Criteria):}
\begin{enumerate}
  \item Trang chi tiết hiển thị các size S, M, L có giá bán; chỉ cho chọn topping đối với món được phép thêm topping.
  \item Giá hiển thị thay đổi theo size và topping đã chọn.
  \item Khi món hoặc topping ngừng bán, giao diện không cho thêm lựa chọn đó vào giỏ.
\end{enumerate}

\subsection{US-05 -- Quản lý giỏ hàng (Task 5)}
\textbf{User Story}\par
Là một khách hàng, tôi muốn thêm, sửa số lượng và xóa món trong giỏ, để kiểm tra các lựa chọn và tổng tiền trước khi tiến hành thanh toán.

\textbf{Tiêu chí chấp nhận (Acceptance Criteria):}
\begin{enumerate}
  \item Món cùng sản phẩm, size và tập topping được gộp vào một dòng; khác size hoặc topping tạo dòng riêng.
  \item Nút tăng, giảm và xóa cập nhật số lượng, tổng tiền và huy hiệu giỏ hàng ngay trên giao diện.
  \item Giỏ đối chiếu sản phẩm, size, topping và định lượng nguyên liệu qua API; khi thiếu một nguyên liệu, khách được báo rõ và không thể tăng quá số ly có thể pha.
  \item Giỏ không bị mất khi tải lại trang; khi khách đăng nhập, các món trong giỏ khách được chuyển vào giỏ tài khoản và gộp đúng theo lựa chọn.
  \item Mã ưu đãi đã chọn được giữ khi tải lại trang và chỉ áp dụng nếu còn hiệu lực, đủ điều kiện.
\end{enumerate}

\subsection{US-06 -- Tạo đơn hàng từ giỏ (Task 6)}
\textbf{User Story}\par
Là một khách hàng đã đăng nhập, tôi muốn gửi giỏ hợp lệ để tạo đơn chờ thanh toán, để nhận mã đơn và chuyển sang trang thanh toán.

\textbf{Tiêu chí chấp nhận (Acceptance Criteria):}
\begin{enumerate}
  \item \texttt{POST /api/orders} yêu cầu đăng nhập và từ chối giỏ rỗng, số lượng sai, món, size hoặc topping không hợp lệ.
  \item Backend tính lại giá từ cơ sở dữ liệu, không tin \texttt{lineTotal} do trình duyệt gửi; mã ưu đãi được kiểm tra lại trước khi lưu.
  \item Backend kiểm tra công thức sản phẩm, topping và lượng nguyên liệu còn dùng được cho toàn bộ giỏ; nếu thiếu, API báo lỗi và không tạo đơn.
  \item Trong một giao dịch cơ sở dữ liệu, hệ thống giữ nguyên liệu và lưu đơn cùng các dòng món. Thành công trả \texttt{orderId}, trạng thái \texttt{PENDING} và thời hạn thanh toán hai phút.
  \item Nếu bất kỳ bước nào thất bại, giao dịch hoàn tác; không còn đơn tạo dở hoặc nguyên liệu bị giữ dở.
\end{enumerate}

\subsection{US-07 -- Đăng ký và đăng nhập (Task 7)}
\textbf{User Story}\par
Là một khách hàng, tôi muốn đăng ký và đăng nhập an toàn, để đặt hàng, xem đơn của mình và nhận điểm thưởng.

\textbf{Tiêu chí chấp nhận (Acceptance Criteria):}
\begin{enumerate}
  \item \texttt{POST /api/auth/register} tạo tài khoản hợp lệ; mật khẩu được băm trước khi lưu.
  \item \texttt{POST /api/auth/login} kiểm tra thông tin đăng nhập và cấp phiên xác thực bằng JWT.
  \item API tạo đơn và xem lịch sử yêu cầu xác thực; một khách không đọc được đơn của khách khác.
\end{enumerate}

\subsection{US-08 -- Thanh toán không tiền mặt (Task 8)}
\textbf{User Story}\par
Là một khách hàng, tôi muốn chọn ví điện tử hoặc thẻ để thanh toán đơn, để hoàn tất việc mua mà không dùng tiền mặt.

\textbf{Tiêu chí chấp nhận (Acceptance Criteria):}
\begin{enumerate}
  \item Trang checkout hiển thị số tiền, mã QR minh họa, phương thức Ví/Thẻ và thời gian còn lại của đơn \texttt{PENDING}.
  \item Khi khách bấm xác nhận, frontend gọi \texttt{POST /api/checkout} hoặc \texttt{POST /api/payments} với \texttt{orderId} và \texttt{Idempotency-Key}; cổng thanh toán trong phạm vi Assignment 1 là mô phỏng.
  \item Thanh toán thành công chuyển đơn sang \texttt{PAID} và trừ nguyên liệu đã giữ; thanh toán lỗi chuyển sang \texttt{PAYMENT\_FAILED} để khách thử lại trong thời hạn.
  \item Đơn quá hạn không thể thanh toán tiếp; nguyên liệu đã giữ được trả về lượng có thể bán.
\end{enumerate}

\subsection{US-09 -- Xác nhận, lịch sử đơn và bàn giao (Task 9)}
\textbf{User Story}\par
Là một khách hàng, tôi muốn xem xác nhận và lịch sử đơn của mình, để kiểm tra khoản thanh toán và theo dõi tiến độ pha chế, nhận món.

\textbf{Tiêu chí chấp nhận (Acceptance Criteria):}
\begin{enumerate}
  \item Sau thanh toán thành công, màn hình hóa đơn hiển thị mã đơn, các món, tổng tiền, phương thức và trạng thái thanh toán.
  \item \texttt{GET /api/orders/me} chỉ trả đơn của tài khoản đang đăng nhập; khách có thể mở chi tiết từng đơn.
  \item Trang lịch sử hiển thị trạng thái từ chờ thanh toán, đã thanh toán, đang pha chế, sẵn sàng đến hoàn tất; trường hợp chưa có đơn có thông báo phù hợp.
  \item Bản bàn giao có Docker Compose cho frontend, backend, PostgreSQL, README và hướng dẫn demo luồng từ chọn món đến xác nhận.
\end{enumerate}

\subsection{US-10 -- Nghiệp vụ backend nâng cao (Task 10)}
\textbf{User Story}\par
Là một quản lý cửa hàng, tôi muốn kiểm soát nguyên liệu, trạng thái đơn, thanh toán lặp và quyền lợi khách hàng, để hệ thống vận hành nhất quán khi nhiều người mua cùng lúc.

\textbf{Tiêu chí chấp nhận (Acceptance Criteria):}
\begin{enumerate}
  \item Chỉ các chuyển trạng thái hợp lệ được chấp nhận theo quy trình \texttt{PENDING} $\rightarrow$ \texttt{PAID} $\rightarrow$ \texttt{PREPARING} $\rightarrow$ \texttt{READY} $\rightarrow$ \texttt{COMPLETED}; trạng thái kết thúc không thể chuyển tiếp tùy ý.
  \item Gửi lại thanh toán với cùng \texttt{Idempotency-Key} không tạo thêm bản ghi thanh toán, không trừ kho hoặc cộng điểm lần thứ hai.
  \item Các đơn đặt đồng thời không được giữ vượt lượng nguyên liệu còn lại; hệ thống kiểm tra định lượng theo công thức sản phẩm/size và topping trong giao dịch.
  \item Mã khuyến mãi chỉ được áp dụng khi còn hiệu lực và đủ điều kiện; điểm thưởng được cộng một lần sau khi thanh toán thành công.
  \item Có kiểm thử chứng minh việc chặn trạng thái sai, chống thanh toán lặp và không bán vượt nguyên liệu khi đặt đồng thời.
\end{enumerate}
